\clearpage\SourceLeaf{146}%
\pdfbookmark[1]{Gauss an von Lindenau}{lindenau1808nov}
\fancyhead[L]{\small GAUSS AN VON LINDENAU}
\begin{center}
{\small GAUSS AN VON LINDENAU.\par}
\medskip\hrule\medskip
{\small Monatliche Correspondenz zur Beförderung der Erd- und Himmels-Kunde,\par
herausgegeben vom Freiherrn \textsc{von Zach}. Januar 1809.\par}
\medskip\hrule
\end{center}\medskip
\begin{flushright}Göttingen, den 30. November 1808.\end{flushright}
Für die mir von Ew. Hochwohlgeb. gütigst mitgetheilten Beobachtungen der Ceres und Juno statte ich Ihnen den verbindlichsten Dank ab. Auch für die Marseiller Cometen-Beobachtungen bin ich Ihnen verbunden, obwohl in denselben \emph{sehr} grosse Fehler begangen zu sein scheinen, und sich daher aus denselben nicht leicht etwas schliessen lassen wird. Heute ist meine Absicht, Ihnen noch einige Anmerkungen über das Problem mitzutheilen, worüber Sie meinen Aufsatz in das October-Heft der \emph{Monatl. Corresp.} aufgenommen haben. \emph{Dieses} Problem ist es eigentlich nicht, was ich \emph{reisenden} Beobachtern vorzüglich empfehlen möchte, sondern mehr dasjenige, welches ich im Anfange jenes Aufsatzes erwähnt und in einem vor kurzem hier gedruckten Programm behandelt habe, wovon ich ein Exemplar beizulegen mir das Vergnügen mache. Das in der \emph{Monatl. Corresp.} abgehandelte Problem wird vorzüglich für solche Beobachter brauchbar sein, die ihren Beobachtungs-Ort mit einem nicht ganz vollkommnen Werkzeuge möglichst scharf bestimmen und zugleich die Fehler des letztern ausmitteln wollen. In diesem Falle hat es gar keine Schwierigkeit, leicht für jede Stunde eine Menge brauchbarer Sterne auszumitteln, wenn man sich voraus für die kenntlichsten Sterne eine Tabelle für die Höhen berechnet, die auch sonst nützlich sein und sehr bequem in die Gestalt einer Karte gebracht werden kann. Man hat gar nicht nöthig sich auf Sterne erster Grösse einzuschränken. Sie se%
\SourceLeaf{147}%
hen aus meinem Beispiele, dass Sterne zweiter Grösse sich auch anwenden lassen, und bei einiger Übung kann man, wenn kein Mondschein ist, sogar Sterne dritter und vierter Grösse noch füglich beobachten, obwohl man dann freilich mehr Sorgfalt nöthig hat, \emph{Verwechslungen} zu vermeiden. Was nun meine Auflösung selbst betrifft, so ist es diejenige, auf die ich sofort von selbst verfiel, als ich ein zur wirklichen Ausübung möglichst bequemes Verfahren suchte, und es fiel mir nicht ein anderswo Auflösungen einer Aufgabe zu suchen, mit der meines Wissens sich noch Niemand beschäftigt hatte. Erst als ich meinen Aufsatz schon abgesandt hatte, fiel es mir auf, dass ein anderes dem Zwecke nach zwar sehr verschiedenes Problem, doch in Ansehung der Auflösung im Wesentlichen mit jenem ganz einerlei ist, das nemlich, wo aus drei \emph{heliocentrischen} Örtern eines Sonnenfleckens die Lage des Sonnen-Äquators und zugleich die Declination des Fleckens gesucht wird; letztere entspricht dann in unserer Aufgabe der Höhe $h$. Mit diesem Problem haben sich bekanntlich eine grosse Menge Geometer beschäftigt, unter denen indessen Niemand eine so zierliche Auflösung gegeben hat, wie \textsc{Cagnoli}. \textsc{La Lande} führt dieselbe in seiner Astronomie an. In Ansehung der Eleganz ziehe ich diese Auflösung der meinigen vor, obwohl ich glaube, dass die Auflösung selbst aus den drei Fundamentalgleichungen zierlicher abgeleitet werden kann, als \textsc{Cagnoli} sie entwickelt hat. Auch in Ansehung der practischen Bequemlichkeit stehen alle übrigen von \textsc{La Lande} angeführten Auflösungen der von \textsc{Cagnoli} weit nach. Die meinige würde noch etwas bequemer sein, wenn man blos $\varphi$ und $k$ sucht, hingegen würde die von \textsc{Cagnoli} etwas kürzer sein, wenn man auch $h$ mit verlangt; im ersten Falle brauche ich 18, \textsc{Cagnoli} 21, im zweiten ich 24, \textsc{Cagnoli} wieder nur 21 Logarithmen; doch ist zu bemerken, dass bei \textsc{Cagnoli} alle 21 zu fast eben so vielen, nemlich zu 19 \emph{verschiedenen} Winkeln gehören, da hingegen bei meiner Auflösung von vielen Winkeln, Sinus und Cosinus oder andere trigonometrische Functionen zugleich aufgesucht werden (so dass, wenn blos $\varphi$ und $k$ gesucht werden, an 14, und wenn auch $h$ mit verlangt wird, an 18 verschiedenen Stellen der Tafeln aufgeschlagen wird), welches allerdings einen Unterschied macht, zumal wenn man sich der schönen \textsc{Taylor}schen Tafeln bedient, wo die ganze Arbeit sich fast blos auf das Aufschlagen reducirt. Die Formeln für unsere Aufgabe nach \textsc{Cagnolis} Auflösung waren folgende:

Man berechne drei Hülfswinkel $A,A',A''$ durch die Gleichungen
\SourceLeaf{148}%
\[
\begin{aligned}
\operatorname{tang}A&=\frac{\sin\frac12(\delta''-\delta')}{\cos\frac12(\delta''+\delta')}\operatorname{cotang}\tfrac12(t''-t')\\
\operatorname{tang}A'&=\frac{\sin\frac12(\delta-\delta'')}{\cos\frac12(\delta+\delta'')}\operatorname{cotang}\tfrac12(t-t'')\\
\operatorname{tang}A''&=\frac{\sin\frac12(\delta'-\delta)}{\cos\frac12(\delta'+\delta)}\operatorname{cotang}\tfrac12(t'-t)
\end{aligned}
\]
mache sodann
\[
\operatorname{tang}B=\frac{\sin\frac12(\delta'-\delta)}{\cos\frac12(\delta'+\delta)}\operatorname{cotang}(A-A')
\]
so ist
\[
k=\tfrac12(t'+t)-B
\]
Man setze ferner $A'+A''-A=C$ und mache
\[
\begin{aligned}
\operatorname{tang}D&=\frac{\cos\frac12(t-k+C)}{\cos\frac12(t-k-C)}\operatorname{tang}(45^0+\tfrac12\delta)\\
\operatorname{tang}E&=\frac{\sin\frac12(t-k-C)}{\sin\frac12(t-k+C)}\operatorname{cotg}(45^0+\tfrac12\delta)
\end{aligned}
\]
\expandafter\def\csname HistoricalNoteText6\endcsname{[Handschriftliche Bemerkung:]
\[
\begin{aligned}
\frac{\cos\varphi}{\cos h}\sin(t-k)&=\sin(A'+A''-A)\\
\frac{\cos\varphi}{\cos h}\sin(t'-k)&=\sin(A''+A-A')\\
\frac{\cos\varphi}{\cos h}\sin(t''-k)&=\sin(A+A'-A'')
\end{aligned}
\]}
so ist\HistoricalNote{6}
\[
\varphi=D+E,\qquad h=D-E
\]

Eine Unbequemlichkeit bei dieser Methode ist, dass man nicht bequem im voraus eine bestimmte Regel geben kann (in so fern man blos nach den analytischen Formeln ohne eine Figur rechnet, ohne von den Beobachtungen etwas weiter als die Uhrzeiten zu entlehnen, also ohne im voraus zu wissen, auf welchen Seiten des Meridians die drei Sterne beobachtet sind), in welchem Halbkreise man die Winkel $A,A',A'',B,D,E$ nehmen müsse, welches bekanntlich die Bestimmung durch die Tangenten für sich unentschieden lässt. Indessen lässt sich zeigen, dass man hiebei einstweilen willkürlich verfahren darf, nur wird man dann in einigen Fällen noch folgende Änderungen machen müssen.

1) Muss man statt $k$, $k+180^0$ oder, welches hier einerlei ist, $k-180^0$ setzen, wenn man für $\varphi$ und $h$ solche Werthe erhalten hat, dass $\cos\varphi$ und $\sin h$ mit entgegengesetzten Zeichen afficirt sind.
\SourceLeaf{149}%

2) Statt $h$ und $\varphi$, wenn man dafür ausserhalb der Grenzen $0$ und $90^0$ liegende Werthe finden sollte, setzt man ihren Unterschied von dem zunächst liegenden Vielfachen von $180^0$.

3) Die Polhöhe ist als nördlich oder südlich zu betrachten, je nachdem $\sin\varphi$ und $\sin h$ gleiche oder entgegengesetzte Zeichen erhalten haben.

Gäben also z.\ B. obige Formeln $\varphi=231^0$, $h=-127^0$, so wäre die wirkliche Polhöhe $51^0$ und die wirkliche Höhe $53^0$. Der Werth von $k$ hingegen bliebe ungeändert, hingegen hätte man gefunden $\varphi=231^0$, $h=+127^0$, so musste $k$ um $180^0$ vermehrt oder vermindert werden. Begreiflich ist dies blos der analytischen Vollständigkeit wegen bemerkt, denn wenn $k$ einer Änderung von $180^0$ bedarf, so ist dies ohnehin klar, da man beim Stande der Uhr um 12 Stunden nicht ungewiss ist.

Der Vergleichung wegen lege ich noch den numerischen Calcul für dasselbe Beispiel nach obigen Formeln bei.

\clearpage\begin{center}\emph{Berechnung des Beispiels nach Cagnolis Methode.}\end{center}
\begin{center}\small
\begin{minipage}[t]{.49\textwidth}\centering
\begin{tabular}[t]{@{}lr@{}}
$\tfrac12(\delta''-\delta')=$&$-24^0\ 49'\ 59''55$\\
$\tfrac12(\delta-\delta'')=$&$-5\quad17\quad25.90$\\
$\tfrac12(\delta'-\delta)=$&$+30\quad7\quad25.45$\\[10pt]
$\tfrac12(t''-t')=$&$-129\quad41\quad39.45$\\
$\tfrac12(t-t'')=$&$+135\quad0\quad4.72$\\
$\tfrac12(t'-t)=$&$-5\quad18\quad25.27$
\end{tabular}
\end{minipage}\hfill
\begin{minipage}[t]{.49\textwidth}\centering
\begin{tabular}[t]{@{}lr@{}}
$\tfrac12(\delta''+\delta')=$&$63^0\ 27'\ 6''15$\\
$\tfrac12(\delta+\delta'')=$&$33\quad19\quad40.70$\\
$\tfrac12(\delta'+\delta)=$&$58\quad9\quad40.25$
\end{tabular}
\end{minipage}
\par\medskip
\begin{tabular}{@{}lr@{}}
$\log\sin\tfrac12(\delta''-\delta')\;.\;.\;.$&$9.6232267\,\mathrm n$\\
$\mathrm{Compl.}\log\cos\tfrac12(\delta''+\delta')\;.$&$0.3497391$\\
$\log\operatorname{cotang}\tfrac12(t''-t)\;.\;.\;.$&$9.9191030$\\\hline
$\log\operatorname{tang}A\;.\;.\;.\;.\;.\;.$&$9.8920688\,\mathrm n$
\end{tabular}
\qquad $A=142^0\quad2\quad50''70$
\par\medskip
\begin{tabular}{@{}lr@{}}
$\log\sin\tfrac12(\delta-\delta'')\;.\;.\;.\;.$&$8.9647590\,\mathrm n$\\
$\mathrm{Compl.}\log\cos\tfrac12(\delta+\delta'')\;.$&$0.0780332$\\
$\log\operatorname{cotang}\tfrac12(t-t')\;.\;.\;.$&$0.0000199\,\mathrm n$\\\hline
$\log\operatorname{tang}A'\;.\;.\;.\;.\;.$&$9.0428121$
\end{tabular}
\qquad $A'=6^0\ 17'\ 51''34$
\end{center}
\SourceLeaf{150}%
\begin{center}\small
\begin{tabular}{lr}
$\log\sin\tfrac12(\delta'-\delta)\;.\;.\;.\;.$&$9.7005892$\\
$\mathrm{Compl.}\log\cos\tfrac12(\delta'+\delta)\;.$&$0.2777516$\\\hline
$(*)\;.\;.\;.\;.\;.\;.\;.\;.$&$9.9783408$\\
$\log\operatorname{cotang}\tfrac12(t'-t)\;.\;.\;.$&$1.0320274$\\\hline
$\log\operatorname{tang}A''\;.\;.\;.\;.\;.$&$1.0103682\,\mathrm n$
\end{tabular}
\par\smallskip
\begin{tabular}{rl}
$A''=$&$95^0\ 34'\ 36''24$\\
$A-A'=$&$135\quad44\quad59.36$\\
$C=$&$-40\quad10\quad23.12$
\end{tabular}
\par\medskip
\begin{tabular}{lr}
$(*)\;.\;.\;.\;.\;.\;.\;.\;.$&$9.9783408$\\
$\log\operatorname{cotang}(A-A')\;.\;.\;.$&$0.0113684\,\mathrm n$\\\hline
$\log\operatorname{tang}B\;.\;.\;.\;.\;.$&$9.9897092\,\mathrm n$
\end{tabular}
\par\smallskip
\begin{tabular}{rl}
$B=$&$315^0\ 40'\ 43''55$\\
$\tfrac12(t+t')=$&$318\quad24\quad44.78$\\\hline
$k=$&$2\quad44\quad1.23$\\[4pt]
$\tfrac12(t-k+C)=$&$140\quad24\quad22.85$\\
$\tfrac12(t-k-C)=$&$180\quad34\quad45.97$\\
$45^0+\tfrac12\delta=$&$59\quad1\quad7.40$
\end{tabular}
\par\medskip
\begin{tabular}{lr}
$\log\cos\tfrac12(t-k+C)\;.\;.$&$9.8868199\,\mathrm n$\\
$\mathrm{Compl.}\log\cos\tfrac12(t-k-C)$&$0.0000222\,\mathrm n$\\
$\log\operatorname{tang}(45^0+\tfrac12\delta)\;.\;.\;.$&$0.2215478$\\\hline
$\log\operatorname{tang}D\;.\;.\;.\;.\;.$&$0.1083899$
\end{tabular}
\par\smallskip
$D=52^0\ 4'\ 36''35$
\par\medskip
\begin{minipage}[t]{\textwidth}\centering
\begin{tabular}{lr}
$\log\sin\tfrac12(t-k-C)\;.\;.$&$8.0048736\,\mathrm n$\\
$\mathrm{Compl.}\log\sin\tfrac12(t-k+C)$&$0.1956297$\\
$\log\operatorname{cotang}(45^0+\tfrac12\delta)\;.\;.$&$9.7784522\,\mathrm n$\\\hline
$\log\operatorname{tang}E\;.\;.\;.\;.\;.$&$7.9789555\,\mathrm n$
\end{tabular}
\par\smallskip
\begin{tabular}{rl}
$E=$&$-0^0\ 32'\ 45''02$\\
Also $\varphi=$&$51\quad31\quad51.33$\\
$h=$&$52\quad37\quad21.37$
\end{tabular}
\end{minipage}
\end{center}
